\documentclass[11pt]{article}
\usepackage{amsmath,amssymb,amsthm}
\usepackage{geometry}
\usepackage{booktabs}
\usepackage{graphicx}
\usepackage{hyperref}
\usepackage{enumerate}
\usepackage{listings}
\usepackage{xcolor}
\usepackage{bm}
\usepackage{fancyhdr}

\geometry{letterpaper, margin=1in}

\lstset{
  basicstyle=\ttfamily\small,
  keywordstyle=\color{blue},
  commentstyle=\color{gray},
  stringstyle=\color{red},
  breaklines=true,
  columns=fullflexible,
  keepspaces=true,
  language=C++
}

\pagestyle{fancy}
\fancyhf{}
\fancyhead[L]{\small\textit{Environmental Particle Extension}}
\fancyhead[R]{\small\textit{VSEPR-SIM v3.0.1}}
\fancyfoot[C]{\thepage}
\renewcommand{\headrulewidth}{0.4pt}

\newtheorem{definition}{Definition}[section]
\newtheorem{remark}{Remark}[section]

\title{\textbf{Environmental Particle Extension} \\[6pt]
\Large Sun/Wind Dual-Kind Carriers, Root Chaos, and Leaf Generation \\[6pt]
\large VSEPR-SIM Development Team}
\author{Formation Engine Research Division}
\date{Version 3.0.1 --- \today}

\begin{document}
\maketitle
\thispagestyle{fancy}

\begin{abstract}
We extend the VSEPR-SIM bead-based framework with a dual environmental
particle regime consisting of radiative \emph{sun particles} and advective
\emph{wind particles}.  Sun particles deposit directional energy into
exposed beads through projected-area and shading-dependent transfer
kernels.  Wind particles contribute momentum loading, drying, and
transport bias.  Root development is modulated by a bounded stochastic
response factor $\Gamma_i \in [0.98,\,1.8]$ combined with piecewise
polynomial moisture response and logic-gated growth modifiers.  Leaf
generation follows a Heaviside step-function initiation gate with
post-threshold continuous expansion.  The system composes additively
with the existing wind-particle perturbation infrastructure.
\end{abstract}

\tableofcontents
\newpage

% =====================================================================
\section{Environmental Particle State Vector}

\begin{definition}[Environmental particle]
Each environmental particle is a lightweight carrier state:
\begin{equation}
  P_k = (\mathrm{id},\;\mathrm{kind},\;\mathbf{r}_k,\;\mathbf{v}_k,\;
         E_k,\;I_k,\;\omega_k,\;\chi_k,\;\tau_k)
\end{equation}
where:
\begin{itemize}
  \item $\mathbf{r}_k$: position
  \item $\mathbf{v}_k$: direction / transport velocity
  \item $E_k$: energy content
  \item $I_k$: local intensity weight
  \item $\omega_k$: oscillation or temporal modulation term
  \item $\chi_k$: stochastic or chaos control parameter
  \item $\tau_k$: lifetime / persistence
\end{itemize}
\end{definition}

The kind field partitions particles into two classes:
\begin{equation}
  \mathrm{kind} \in \{\,\textsc{Sun},\;\textsc{Wind}\,\}
\end{equation}

Sun particles are not literal photons.  They are radiative packets or
directional energy carriers.  Wind particles are not fluid elements.
They are momentum plus transport disturbance packets.

% =====================================================================
\section{Sun Particle Deposition}

For a bead $i$, the energy deposited by sun particle $k$:
\begin{equation}
  \Delta E_i^{\mathrm{sun}} = \sum_{k \in \mathrm{SUN}}
    A_i^{\mathrm{proj}}\;I_k\;T_i\;S_i^{\mathrm{photo}}\;
    \Phi_{\mathrm{shade}}(i, k)
  \label{eq:sun}
\end{equation}

\begin{center}
\begin{tabular}{ll}
\toprule
\textbf{Symbol} & \textbf{Meaning} \\
\midrule
$A_i^{\mathrm{proj}}$ & Projected area toward incident radiation \\
$T_i$ & Tissue transmissivity or absorptivity \\
$S_i^{\mathrm{photo}}$ & Biological photo-response coefficient \\
$\Phi_{\mathrm{shade}}(i,k)$ & Shadowing / occlusion factor $\in [0,1]$ \\
\bottomrule
\end{tabular}
\end{center}

For leaves, sun deposition drives:
\begin{itemize}
  \item Growth activation
  \item Local expansion bias
  \item Surface heating
  \item Transpiration tendency
\end{itemize}

For roots, sun is indirect and arrives through temperature field
coupling, moisture gradient response, and shoot-driven resource
allocation.

% =====================================================================
\section{Wind Particle Force}

\begin{equation}
  \Delta\mathbf{F}_i^{\mathrm{wind}} = \sum_{k \in \mathrm{WIND}}
    C_{d,i}\;A_i^{\mathrm{eff}}\;\|\mathbf{v}_k - \mathbf{v}_i\|^2\;
    \hat{\mathbf{n}}_{ik}\;\Psi_i^{\mathrm{flex}}
  \label{eq:wind}
\end{equation}

\begin{center}
\begin{tabular}{ll}
\toprule
\textbf{Symbol} & \textbf{Meaning} \\
\midrule
$C_{d,i}$ & Drag-like bead coefficient \\
$A_i^{\mathrm{eff}}$ & Effective exposed area \\
$\hat{\mathbf{n}}_{ik}$ & Loading direction (unit relative velocity) \\
$\Psi_i^{\mathrm{flex}}$ & Structural flexibility response \\
\bottomrule
\end{tabular}
\end{center}

Wind drives not only force but also:
mechanical bending, oscillatory fatigue, pollen transport, evaporation
and drying, boundary-layer disturbance, seed release probability.

The drying rate from wind:
\begin{equation}
  \dot{D}_i = \kappa_{\mathrm{dry}}\;\|\mathbf{v}_k - \mathbf{v}_i\|\;A_i^{\mathrm{eff}}
\end{equation}

% =====================================================================
\section{Root Chaos Factor}

\subsection{Bounded Stochastic Multiplier}

The root exploration multiplier:
\begin{equation}
  \Gamma_i = \mathrm{clamp}\!\bigl(
    \alpha_0 + \alpha_1 M_i + \alpha_2 \nabla W_i
    + \alpha_3 \rho_i^{\mathrm{soil}} + \alpha_4 \chi_i,
    \;0.98,\;1.8
  \bigr)
  \label{eq:gamma}
\end{equation}

\begin{center}
\begin{tabular}{lll}
\toprule
\textbf{Symbol} & \textbf{Meaning} & \textbf{Default} \\
\midrule
$\alpha_0$ & Base offset & 1.0 \\
$\alpha_1$ & Moisture weight & 0.3 \\
$\alpha_2$ & Water gradient weight & 0.15 \\
$\alpha_3$ & Soil density penalty & $-0.2$ \\
$\alpha_4$ & Stochastic amplitude & 0.1 \\
\bottomrule
\end{tabular}
\end{center}

This gives:
\begin{itemize}
  \item Near-deterministic roots in stable media
  \item More chaotic probing in poor or heterogeneous zones
  \item Occasional lateral overreach and branching competition
\end{itemize}

\subsection{Piecewise Polynomial Response}

\begin{equation}
  R_{\mathrm{root}}(x) =
  \begin{cases}
    a_1 x^2 + b_1 x + c_1 & x < x_{\mathrm{dry}} \\
    a_2 x^3 + b_2 x^2 + c_2 x + d_2 & x_{\mathrm{dry}} \leq x < x_{\mathrm{opt}} \\
    a_3 x^2 + b_3 x + c_3 & x \geq x_{\mathrm{opt}}
  \end{cases}
\end{equation}

The piecewise polynomial is wrapped in logic:
\begin{equation}
  g_i^{\mathrm{root}} = \mathrm{clamp}\!\bigl(
    R_{\mathrm{root}}(M_i) \cdot L_{\mathrm{gates}}(s_i)
    \cdot \Gamma_i,\;0.05,\;2.25
  \bigr)
\end{equation}

where $L_{\mathrm{gates}}$ applies multiplicative penalties:
\begin{align}
  \text{if } D_i > D_{\max}: &\quad \times\, 0.55 \notag\\
  \text{if } \rho_i^c > \rho_{\max}: &\quad \times\, 0.72 \notag\\
  \text{if } \nabla N_i > N_{\mathrm{trig}}: &\quad \times\, 1.18 \notag\\
  \text{if } S_i^{\mathrm{sun}} < S_{\min}: &\quad \times\, 0.91
\end{align}

% =====================================================================
\section{Leaf Generation Gate}

\subsection{Heaviside Initiation}

\begin{equation}
  H_{\mathrm{leaf}} = H\!\bigl(
    \beta_1 E_{\mathrm{local}} + \beta_2 S_{\mathrm{light}}
    + \beta_3 T_{\mathrm{hydr}} - \beta_4 \Sigma_{\mathrm{damage}}
    - \Theta_{\mathrm{leaf}}
  \bigr)
\end{equation}

Below threshold: no leaf initiation.  Above threshold: leaf primordium
appears.

\subsection{Continuous Expansion}

After initiation, growth becomes continuous:
\begin{equation}
  \frac{dA_{\mathrm{leaf}}}{dt} = H_{\mathrm{leaf}} \cdot
  G_{\mathrm{leaf}}(E, L, W, \Sigma)
\end{equation}

where $G_{\mathrm{leaf}}$ depends on energy availability, light, water,
and structural stress.

% =====================================================================
\section{Energy Decomposition}

The total environmental energy:
\begin{equation}
  U_{\mathrm{env}} = U_{\mathrm{wind}} + U_{\mathrm{sun}}
    + U_{\mathrm{dry}} + U_{\mathrm{photo}} + U_{\mathrm{stress}}
\end{equation}

Each term is explicitly tracked and inspectable, consistent with the
anti-black-box design philosophy.

% =====================================================================
\section{Particle Advection and Lifetime}

Particles advance by:
\begin{align}
  \mathbf{r}_k(t + \Delta t) &= \mathbf{r}_k(t) + \mathbf{v}_k\,\Delta t \\
  \tau_k(t + \Delta t) &= \tau_k(t) - \Delta t
\end{align}

With optional oscillation modulation:
\begin{equation}
  I_k \leftarrow I_k \cdot \bigl(1 + 0.1\sin(\omega_k\,\Delta t)\bigr)
\end{equation}

A particle is removed when $\tau_k \leq 0$.

% =====================================================================
\section{Integration with Existing Wind Particle System}

The existing \texttt{WindParticle} (from \texttt{wind\_particle.hpp})
operates on the atomistic \texttt{State} directly.  The environmental
particle extension operates at the bead level.  The two compose:

\begin{enumerate}
  \item \texttt{WindParticle::apply(state)} adds force to atomistic state
  \item \texttt{interact\_env\_particle(bead, p)} computes plant-bead response
  \item Both contribute to $U_{\mathrm{ext}}$ additively
\end{enumerate}

The same codebase serves both metallic/turbine and plant workflows.

\end{document}
